\documentclass[hyperref,UTF8,12px,a4paper]{ctexart} % use larger type; default would be 10pt
%a4paper: width = 17cm

% \usepackage[utf8]{inputenc} % set input encoding (not needed with XeLaTeX)

%%% PAGE DIMENSIONS
\usepackage[margin=1in]{geometry} % to change the page dimensions
%\geometry{a4paper} or letterpaper (US) or a5paper or....
% \geometry{margin=2in} % for example, change the margins to 2 inches all round
% \geometry{landscape} % set up the page for landscape
%   read geometry.pdf for detailed page layout information
\usepackage{xeCJK}
\usepackage{graphicx} % support the \includegraphics command and options

% \usepackage[parfill]{parskip} % Activate to begin paragraphs with an empty line rather than an indent

%%% PACKAGES
\usepackage{abstract} % for abstract
\usepackage{amsmath}
\usepackage{booktabs} % for much better looking tables
\usepackage{cite} % reference
\usepackage{array} % for better arrays (eg matrices) in maths
\usepackage{paralist} % very flexible & customisable lists (eg. enumerate/itemize, etc.)
\usepackage{verbatim} % adds environment for commenting out blocks of text & for better verbatim
\usepackage{subfigure} % make it possible to include more than one captioned figure/table in a single float
% These packages are all incorporated in the memoir class to one degree or another...
\usepackage{cases} %equation set
\usepackage{multirow} %use table
\usepackage{lastpage} %get the total page
\usepackage{algorithm}
\usepackage{algorithmic}
\usepackage{enumerate}
\usepackage{wrapfig}
\usepackage{comment} % comment for more than one line
\usepackage{hyperref}
\hypersetup{colorlinks,linkcolor=black,anchorcolor=black,citecolor=black,
pdfstartview=FitH,bookmarksnumbered=true,bookmarksopen=true,} % set href in tex & pdf

%%% HEADERS & FOOTERS
\usepackage{fancyhdr} % This should be set AFTER setting up the page geometry
\pagestyle{fancy} % options: empty , plain , fancy
\fancyhf{}
\renewcommand{\headrulewidth}{0.5pt} % customise the layout...
\renewcommand{\footrulewidth}{0.5pt}
\lhead{}\chead{学号：xxx\quad 姓名：yyy}\rhead{}
\lfoot{}\cfoot{第 \thepage 页，共 \pageref{LastPage} 页}\rfoot{}

%%% SECTION TITLE APPEARANCE
\usepackage{sectsty}
\allsectionsfont{\sffamily\mdseries\upshape} % (See the fntguide.pdf for font help)
% (This matches ConTeXt defaults)
\usepackage{indentfirst}
\setlength{\parindent}{2em}
\numberwithin{equation}{section} %formed the equation number as (1.1)...
\renewcommand\contentsname{\heiti 目录}
\renewcommand\refname{\heiti 参考文献}
%%% ToC (table of contents) APPEARANCE
\usepackage[notlot,notlof]{tocbibind} % Put the bibliography in the ToC
\usepackage[titles,subfigure]{tocloft} % Alter the style of the Table of Contents
%\renewcommand{\cftsecfont}{\rmfamily\mdseries\upshape}
%\renewcommand{\cftsecpagefont}{\rmfamily\mdseries\upshape} % No bold!

%%% END Article customizations

%%% The "real" document content comes below...

\title{作业}
\author{first\\
second}
%\date{} Activate to display a given date or no date (if empty),
         % otherwise the current date is printed

\begin{document}
\maketitle
\begin{figure}[!htb]
	\centering
	\includegraphics[width=4.5cm,height=7.5cm]{nkbackground.png}
\end{figure}
\thispagestyle{fancy} % set the page style of the first page

{\noindent \Large \heiti 摘要}\\
% Chinese abstract
\par
{\noindent \Large \textbf{Abstract}}\\
% English abstract
\par
{\noindent \Large \heiti 关键词}\\
% Chinese keywords
\par
{\noindent \Large \textbf{Keywords}}\\
% English keywords
\section{第一章}
\subsection{第一节}
% \tableofcontents
\newpage
\begin{figure}[!htb]
	\centering
    \includegraphics[width=12cm]{nkbackground.png}
    \caption{设计草图 by xxx}
\end{figure}

\begin{description}
    \item[抗摔耐用] 1m抗摔，在精度前提下考虑防护
    \item[两次注胶] 对模具不仅进行初注胶，而且进行重复加工软注胶
    \item[精度更强] 全铝合金传感器，并进行黑体校准
    \item[蜂鸣报警] 超过量程时进行蜂鸣报警
    \item[屏幕显示] 我们可以显示最大值（MAX），最小值（MIN），最多存储数据（STO），高温预警（HAL），低温预警（LAL），发射率设置（EMS），平均值（AVG），基本满足民用需求
\end{description}

\newpage
\begin{figure}[!htb]
	\centering
    \includegraphics[width=15cm]{nkbackground.png}
    \caption{设计草图 by yyy}
\end{figure}

\newpage
\begin{figure}[!htb]
	\centering
    \includegraphics[width=15cm]{nkbackground.png}
    \caption{设计草图 by zzz}
\end{figure}

\begin{comment}
\begin{itemize}
    \item 加入激光指示，远距离自动放大激光点大小
    \item 握柄采用流线型设计
    \item 过一段时间后可以自动关机
\end{itemize}
\end{comment}
\end{document}
